\documentclass[a4paper,11pt]{article}

\usepackage{revy}
\usepackage[utf8]{inputenc}
\usepackage[T1]{fontenc}
\usepackage[danish]{babel}


\revyname{MatematikRevyen}
\revyyear{2026}
\version{1}
\eta{$1.5$ minutter}
\status{Færdig}

\title{Matematik på mellemtrinnet}
\author{Julius Pallesen '24}
\melody{Annika Hoydal: ``Dukka mín er blá''}

\begin{document}
\maketitle

\begin{roles}
\role{S}[Sanger] Sanger evt. med guitar (eller balalajka el.lign.) på scenen
\end{roles}

\begin{song}
    \sings{S}[Vers] Én plus én er to
    To plus to er fire
    Tre plus tre er seks
    Fir' plus fir' er otte
    Fem plus fem er ti

    \sings{S}[Vers] På 0. klassetrin
    Sku' vi lægge sammen
    Så kom mellemtrin
    Der sku' vi lære at gange
    Og lære kvadrattal

    \scene{Kort mellemspil}

    \sings{S}[Vers] Én gang' én er én
    To gang' to er fire
    Tre gang' tre er ni
    Fir' gang' fir' er seksten
    Fem gang' fem er femogtyve
\end{song}

\scene{Note: Denne sang kan evt. fremføres under et sceneskift af en revyst placeret i publikum (med spot), der hiver en guitar frem og fremfører den akustisk}

\end{document}

